\documentclass[USenglish]{uiophdthesis}
\usepackage[utf8]{inputenc}
\usepackage[T1]{url}\urlstyle{sf}
\usepackage{babel, csquotes, graphicx, textcomp, varioref}
\usepackage[backend=biber,style=numeric-comp]{biblatex}
\usepackage[hidelinks, hypertexnames=false]{hyperref}

\title{Decoding transcription factor binding to the DNA}
\author{Ieva Rauluševičiūtė} 

\addbibresource{thesisbib.bib}

\begin{document}
\frontmatter{}
\maketitle[
  supervisor={Anthony Mathelier},
  % co-supervisor={Jonas Paulsen}, %% if more than one
  %% dept={Department of Something},  %% if you want to include it
  %% fac={Faculty of Theory},         %% if you want to include it
]

\begin{abstract}
  TODO: Abstract of the thesis.
\end{abstract}

\begin{xabstract}[Sammendrag]
  TODO: abstract in Norwegian (mandatory)
\end{xabstract}

\tableofcontents{}
\listoffigures{}
\listoftables{}

\begin{preface}
  TODO: preface (acknowledgements and thanks)
\end{preface}

\mainmatter{}






\part*{Introduction}
  \chapter{Background}
    \section{Organisation of the genome}

    Talk about Eukaryotic genomes only...

    A eukaryotic genome contains different types of repetetive regions or motifs. These include transposable elements, tandem repeats, GC- and AT-rich sequences, splice sites and binding sites of non-coding RNA molecules. Importantly, crutial for gene expression regulation are DNA sequence motifs targeted by transcription factors \cite{boeva2016analysis}.


      \subsection{DNA}

      DNA existence was discovered in 1869 by Friedrich Miescher when he found a new substance inside the nuclei of the human white blood cells that did not have the same properties as proteins or lipids \cite{pray2008discovery,dahm2008discovering}.
              
      However, science world often emphasize not the discovery of nucleic acid itself, but the publishing of the structure of DNA. This was a fundamental moment in science as it allowed to explain the possible mechanism of DNA replication, thus information multiplication and transfer.



      Decipheric DNA structure
      This provided an explanation on how proteins are encoded by genes - nucleotides in the DNA encode the amino acids in the protein chain





      \subsection{Chromatin}






    \section{Fundamentals of gene transcription regulation}




    Transcriptional regulation - at the DNA level:
    - chromatin state modification (histone tail post-translational modifications - histone marks)
        Histone methyltransferases, histone acetyltransferases, histone deacetylates, polycomb proteins.
    - DNA methylation
        DNA methyltransferases, 5mC-reading proteins (could be a lead to the transcription factors, slightly mention the poineering function, but more later)
    - TFs
        Bind enhancers, making a loop. (mention this very briefly, more details later on)



    Post-transcriptional regulation - at the RNA level:
    - small RNAs
    - capping, RNA splicing, poly(A) tail, RNA editing, mRNA stability, nuclear export.




    (Focus on transcriptional regulation.)





    Gene expression regulation by transcription factors is one of the most studied fields, yet still contains a number of unanswered questions that keep the scientists busy. <---- as a last sentence?

  \chapter{Gene transcription regulation by transcription factors}




    \section{Transcription factors}





      \subsection{Protein structures and classification}







        \subsubsection{Protein domains}






        \subsubsection{Classification}

        DNA binding domains are highly conserved, thus are used to classify these proteins into families and classes.

        One of the most extensive classification of mammalian transcription factors - TFClass - include superclass, class, family, subfamily and genus \cite{wingender2018tfclass, wingender2013criteria}.

    \section{Transcription factor binding}


    Transcription factors directly or indirectly bind genomic regions, known as promoters and enhancers. The location where TF resides onto DNA is referred to as a transcription factor binding site (TFBS). A promoter or an enhancer can have multiple TFBSs of the same or different TF. Transcription factors generally have binding preference to a set of DNA sequences - DNA binding motif \cite{boeva2016analysis}.



      \subsection{DNA binding models}

      In order to interpret TF binding to the regulatory regions, we need ways to represent the DNA binding. Currently, DNA binding is represented using various models.
      
      The most popular and widely used is positional weight matrix (PWM) or position-specific scoring matrix (PSSM). 


      PFM, PWM. Other models to represent binding.


      


      There are two main strategies applied to analyse genomic sequences from TF binding perspective. One is to find over-represented sequence motifs using existing collections of DNA binding motifs. The second is to discover over-represented motifs \textit{de novo} \cite{boeva2016analysis}. 


      \subsection{Mechanisms of gene regulation}


      Go from the simplest to more complicated....

        \subsubsection{Pioneering}

        TFs are capable of binding heterochramatin through recruitment of histone acetyltransferases \cite{boeva2016analysis}.


        
        \subsubsection{Cooperative transcription factor binding}


        When TFs interact or compete with each other binding sites co-localize or overlap \cite{boeva2016analysis}.

        Notes: talk about mechanisms of binding and methods that investigate these mechanisms can be explained here too.
        So this needs to be quite specific mechanisms and then specifics on the methods as well (experimental and computational methods that is).









  \chapter{Cell type-specific gene expression regulation}

  This is about how gene expression is cell-type specific, relate to the fact that we have some cool single-cell techniques that tackle this problem, but we have a lot of bulk data that could still be used.



    \section{Gene expression regulation and cancer}


    \section{Gene expression data deconvolution}





















\part*{Objectives of the Study}



\part*{Summary of Papers}

\chapter*{Papers I-II}

\section*{Paper I}

\section*{Paper II}

\chapter*{Paper III}

\chapter*{Paper IV}







\part*{Discussion}
% \part*{Conclusion}                     %% ... or ?? (if any)
% \chapter{Results}                     %% ... or ??
% Needed?

% \chapter{Discussion}                     %% ... or ??
% TODO







\backmatter{}
\printbibliography{}




\uiopaper{JASPAR 2022: the 9th release of the open-access database of transcription factor binding profiles}
\uioincludepdf{jaspar_2022_paper.pdf}

\uiopaper{JASPAR 2024: }
%% \uioincludepdf{jaspar_2024_paper.pdf}

\uiopaper{Identification of transcription factor co-binding patterns with non-negative matrix factorization}
\uioincludepdf{cobind_paper.pdf}


\uiopaper{TFmeCan deconvolution}
\chapter{Manuscript}
TODO: manuscript or draft


\end{document}
